\documentclass[10pt,letterpaper]{article}

\usepackage[T1]{fontenc}
\usepackage{mathptmx}
\usepackage[letterpaper,left=0.55in,right=0.55in,top=0.35in,bottom=0.35in]{geometry}
\usepackage{array}
\usepackage{tabularx}
\usepackage[hidelinks]{hyperref}

\input{glyphtounicode}
\pdfgentounicode=1

\hypersetup{
  pdftitle={Chrislyr John P. Tan - Curriculum Vitae},
  pdfauthor={Chrislyr John P. Tan},
  pdfsubject={Curriculum Vitae}
}

\pagestyle{empty}
\setlength{\parindent}{0pt}
\setlength{\parskip}{0pt}
\setlength{\tabcolsep}{0pt}
\renewcommand{\arraystretch}{1}
\linespread{0.97}
\urlstyle{same}

\newenvironment{cvitems}{%
  \begin{list}{\textbullet}{%
    \setlength{\leftmargin}{1.25em}%
    \setlength{\labelsep}{0.45em}%
    \setlength{\itemsep}{0pt}%
    \setlength{\topsep}{1.5pt}%
    \setlength{\parsep}{0pt}%
    \setlength{\partopsep}{0pt}%
  }%
}{\end{list}}

\newcommand{\cvsection}[1]{%
  \vspace{7pt}%
  {\fontsize{12}{13}\selectfont\bfseries\MakeUppercase{#1}}\par
  \vspace{2pt}\hrule height 0.4pt\vspace{4pt}%
}

\newcommand{\headingrow}[2]{%
  \begin{tabularx}{\textwidth}{@{}>{\raggedright\arraybackslash}Xr@{}}
    \textbf{#1} & \textbf{#2}
  \end{tabularx}%
}

\newcommand{\subheadingrow}[2]{%
  \begin{tabularx}{\textwidth}{@{}>{\raggedright\arraybackslash}Xr@{}}
    \textit{#1} & #2
  \end{tabularx}%
}

\newcommand{\projectheading}[3]{%
  \begin{tabularx}{\textwidth}{@{}>{\raggedright\arraybackslash}Xr@{}}
    \textbf{#1} \textbar{} \textit{#2} & \textbf{#3}
  \end{tabularx}%
}

\begin{document}

\begin{center}
  {\fontsize{20}{22}\selectfont\bfseries CHRISLYR JOHN P. TAN}\par
  \vspace{1pt}
  {\fontsize{10}{11}\selectfont
    Cebu City, Philippines \textbar{} +63 932 423 0938 \textbar{}
    \href{mailto:cjtan2406@gmail.com}{cjtan2406@gmail.com} \textbar{}
    \href{https://github.com/choonhows}{github.com/choonhows} \textbar{}
    \href{https://choonhows.github.io}{choonhows.github.io}}
\end{center}
\vspace{-4pt}

\cvsection{Education}
\headingrow{University of San Carlos}{Cebu City, Philippines}
\subheadingrow{Bachelor of Science in Computer Science}{2024--Present}
\begin{cvitems}
  \item Third-year undergraduate student; expected graduation: 2027.
\end{cvitems}
\vspace{1pt}
\headingrow{University of San Jose-Recoletos}{Cebu City, Philippines}
\subheadingrow{Technical-Vocational-Livelihood (TVL), Information and Communication Technology}{2022--2024}

\cvsection{Certification}
\headingrow{PhilNITS Fundamental Information Technology Engineers (FE) Certification Examination}{April 2026}
Passer

\cvsection{Projects}
\projectheading{Caledoro}{Flutter, Dart, Riverpod, Hive, Gemini API}{2026}
\href{https://github.com/choonhows/Caledoro}{github.com/choonhows/Caledoro}
\begin{cvitems}
  \item Collaborated as one of three developers on an offline-first productivity application featuring task management, Pomodoro sessions, calendar tools, and Android home-screen widgets.
  \item Integrated the Gemini API to analyze user-created tasks and generate smaller, actionable subtasks, helping users break down complex tasks into manageable steps.
  \item Contributed to feature implementation, debugging, quality assurance, local data persistence, state management, notifications, and responsive application behavior.
  \item Performed detailed testing and iterative refinement to improve reliability, usability, and overall user experience.
\end{cvitems}
\vspace{2pt}

\projectheading{NitPicker}{React, TypeScript, Vite, Tailwind CSS}{2026}
\href{https://github.com/choonhows/NitPicker}{github.com/choonhows/NitPicker}
\begin{cvitems}
  \item Contributed to a PhilNITS FE preparation platform with configurable mock exams, answer explanations, performance analytics, and previous-exam browsing.
  \item Implemented UI refinements including skeleton loading states and quality-of-life improvements for smoother user interaction.
  \item Conducted functional, interactivity, and usability testing and collaborated on responsiveness and workflow refinements.
\end{cvitems}
\vspace{2pt}

\projectheading{Chronos}{Python, FastAPI, PostgreSQL, React, TypeScript, DEAP, Mlxtend}{2026--Present}
\href{https://github.com/saiimonn/Chronos}{github.com/saiimonn/Chronos}
\begin{cvitems}
  \item Collaborated as one of four developers on a locally hosted automated class scheduling system that uses Association Rule Mining to guide a Genetic Algorithm while keeping hard constraints authoritative.
  \item Implemented FastAPI and SQLAlchemy services and React/TypeScript workflows for authentication and role-based access control, scheduling, validation, persistence, timetable projection, and CSV/XLSX export.
  \item Built unit and PostgreSQL integration tests for migrations, APIs, scheduling algorithms, and persistence.
  \item Designed and operationalized a Jira-GitHub ticketing workflow with dual-ID traceability and AI-assisted MCP commands for ticket intake, branch creation, pull-request review, CI checks, and issue completion to streamline team delivery.
\end{cvitems}

\cvsection{Technical Skills}
\textbf{Programming Languages:} Dart, TypeScript, C\#, Python, SQL\par
\textbf{Frameworks \& Technologies:} Flutter, React, Vite, Tailwind CSS, .NET MAUI, FastAPI, SQLAlchemy, PostgreSQL, DEAP, Mlxtend\par
\textbf{Development Tools:} Git, GitHub, Visual Studio Code, Jira, Docker, Supabase\par
\textbf{Mobile \& Web Development:} Responsive UI development, local data persistence, state management, debugging, quality assurance, REST APIs\par
\textbf{AI-Assisted Development:} OpenCode, npx-based skills workflows, AI prompting, loop engineering, prompt engineering, context management, token-efficient prompting, Jira/GitHub MCP workflows\par
\textbf{Software Quality:} Functional testing, UI testing, debugging, usability review, code refinement, Pytest, Vitest, integration testing

\cvsection{Core Competencies}
Problem Solving \enspace\textbullet\enspace Critical Thinking \enspace\textbullet\enspace Attention to Detail
\enspace\textbullet\enspace Team Collaboration \enspace\textbullet\enspace Adaptability
\enspace\textbullet\enspace Continuous Learning\\
\textbullet\enspace Technical Communication \enspace\textbullet\enspace Time Management

\end{document}
